\section{Introduction}

Neuromorphic intermediate representations make model structure portable, but
execution remains platform-specific. NIR intentionally describes continuous-time
primitives without fixing a digital discretization, and its cross-platform study
shows that differences in neuron definitions, integration, threshold timing,
quantization, and nondeterminism can change deployed dynamics
\citep{pedersen2024nir}. Recurrent connections can amplify a local transition
difference over time. A successful conversion therefore establishes that a graph
can be executed, not that its predictions have been preserved.

This distinction creates a practical decision problem. A deployer has trained
weights, an executable reference, a candidate backend, and unlabeled examples
from the intended distribution. Before paying for a full deployment---and before
revealing evaluation labels---the deployer wants to know whether the target's
absolute accuracy change is below a stated budget. If not, the deployer wants a
repair that changes only permitted neuron or numeric parameters and can itself be
re-certified.

Prediction disagreement is a natural observable, but it is not a new idea.
Disagreement has been analyzed for label-free performance estimation under
distribution shift \citep{lee2023dotl}, and disagreement discrepancy has been used
to derive error bounds \citep{rosenfeld2023disagreement}. Our setting uses a
simpler deterministic inequality: changing correctness on any labeled example
requires changing the prediction. The scientific question is consequently not
whether disagreement correlates with loss. It is whether one can prospectively
bound disagreement for a declared \emph{family of SNN execution semantics},
account for the gap between an emulator and physical hardware, and make the
result tight enough to direct repair.

We pursue four contributions:
\begin{enumerate}
    \item an operational, versioned execution-semantics object that fixes the
    digital choices omitted by a model graph;
    \item a hierarchy from exact mappings through sound per-input family analysis
    to simultaneous distribution and physical certificates;
    \item label-free certificate-directed calibration selected separately from
    the audit used for the final bound; and
    \item prospective validation on SpiNNaker2 and a new parameterized Virtex-7
    SRNN realization under a preregistered two-task, five-seed protocol.
\end{enumerate}

The intended claim is conditional and falsifiable. We will not claim that two
update equations are equivalent because a pilot trace happened to match. We will
not call an emulator-only result a physical certificate. We will not proceed with
the present framing if the static bound is mostly vacuous or physical conformance
dominates the total risk.

